\section{Exact traces over a curve}
\label{sec:curve-fibration-traces}

An exact representative on a curve also gives higher-dimensional examples by pullback. Multiplier ideals commute with pullback by a submersion, so the trace data on a modification is given by the trace data on the curve, multiplied by the multiplicities of the pulled-back fibers. The pulled-back Green potential has precisely these divisorial Lelong numbers, and its domination property identifies the ordinary type of the trace envelope.

In this section, $p:X\to M$ is a surjective holomorphic submersion of connected compact complex manifolds with connected fibers. Let $C\subset M$ be a connected compact smooth curve, put $Y=p^{-1}(C)$, and write $q=p|_Y:Y\to C$; then $Y$ is a connected compact submanifold of $X$. Take $U\in\QPSH(M)$ with $\nu(U,C)=0$, and set $u=U\circ p$. The generic Lelong number of $u$ along $Y$ is zero, since in local product coordinates it equals that of $U$ along $C$. Write $\tau_{C,U}$ for the trace data of $U$ along $C$, defined as in \cref{sec:restricted-ideals} with $(M,C,U)$ in place of $(X,Y,u)$, and let
\[
 D=\sum_{b\in C}\tau_{C,U}(b)\,\delta_b,
 \qquad \delta=D(C)<+\infty.
\]
Choose a smooth positive form $\eta$ of integral one on $C$ and a nonpositive Green kernel for $\eta$. The construction in the proof of \cref{thm:curve-trace-realization} gives $w\in\QPSH(C)$, not identically $-\infty$, with
\[
 \ddc w=D-\delta\eta.
\]
\Needspace{9\baselineskip}
For $b\in C$, put $Y_b=q^{-1}(b)$, and let $[Y_b]$ be its current of integration. Set $v=w\circ q$. Pullback of distributions by a submersion is defined by duality with integration along the fibers; it commutes with differentiation and, in the present situation, with positive sums, so
\begin{equation}\label{eq:curve-fiber-divisorial-current}
 \ddc v=\sum_{b\in C}\tau_{C,U}(b)[Y_b]-\delta q^*\eta.
\end{equation}
Each fiber is a smooth connected prime divisor, and the sum is a positive closed current of finite mass, as also follows locally from the positivity of $\ddc(w\circ q)+\delta q^*\eta$.

\Needspace{24\baselineskip}
\begin{theorem}\label{thm:curve-fibration-exact-trace}\index{Trace!curve fibration}
The potential $v$ realizes the entire divisorial trace data of $u$ along $Y$:
\[
 \nu(v\circ\pi,E)=\tau_{Y,u}(E)
\]
for every prime divisor $E$ on every smooth compact modification $\pi:Z\to Y$. More explicitly, write $h=q\circ\pi$. If $h(E)=C$, both sides are zero. If $h(E)=\{b\}$, choose a local coordinate $t_b$ on $C$ vanishing at $b$; then both sides equal
\[
 \operatorname{ord}_E(t_b\circ h)\,\tau_{C,U}(b).
\]

Choose any Hermitian form $\Omega_0$ with $v\in\PSH(Y,\Omega_0)$, and set $\Omega=2\Omega_0$. Then
\begin{equation}\label{eq:curve-fibration-trace-envelope}
 \mathcal T_\Omega(u;Y)=P_\Omega[v]=P_\Omega[v]_{\mathcal I}\sim v.
\end{equation}
The envelope has positive finite extended mass, and the conclusions of \cref{cor:trace-envelope-prescribed-equation} on equations in its ordinary type hold. Its ordinary type does not depend on the choice of the Green kernel, of the form $\eta$, or of the sufficiently large background $\Omega$.
\end{theorem}
\begin{proof}
\Cref{lem:submersion-multiplier-pullback} and restriction to $Y$ give
\[
 \mathfrak a_t(u;Y)=\mathfrak a_t(U;C)\mathcal O_Y.
\]
The nonzero coherent ideal $\mathfrak a_t(U;C)$ on the smooth curve is generated near $b$ by $t_b^{k_t(b)}$, where $k_t(b)=\operatorname{ord}_b\mathfrak a_t(U;C)$. For a divisor $E$ with $h(E)=\{b\}$, the order along $E$ of the extended ideal is therefore $k_t(b)\operatorname{ord}_E(t_b\circ h)$. For a divisor dominating $C$, the generic order is zero, since the ideal on the curve has only finitely many zeros. Dividing by $t$ and letting $t\to\infty$ gives exactly the stated formula for $\tau_{Y,u}(E)$.

For the Lelong numbers of $v\circ\pi$, suppose first that $h(E)=\{b\}$. At a general point of $E$, a local equation $z_1=0$ of $E$ satisfies $t_b\circ h=z_1^k g$ with $g$ holomorphic and nonvanishing, and $k=\operatorname{ord}_E(t_b\circ h)$. Replacing $z_1$ by $z_1g^{1/k}$ for a local holomorphic $k$th root, we may assume $t_b\circ h=z_1^k$. The Lelong number of the one-variable function $w(z_1^k)$ at $0$ is $k\nu(w,b)=k\tau_{C,U}(b)$, as follows from the definition through circular means, and this is also the generic Lelong number of $v\circ\pi$ along $E$ in these product coordinates. If $h(E)=C$, the regular locus of $E$ contains a point mapped to a point where $w$ is finite, because the nonconstant holomorphic map from that locus to $C$ has open image; the pullback is finite there, so its generic Lelong number along $E$ is zero. This proves exact realization on every compact model.

By \cref{lem:divisorial-potential-domination} and \eqref{eq:curve-fiber-divisorial-current}, a quasi-plurisubharmonic function $a$ on $Y$ satisfies $a\preceq v$ if and only if $\nu(a,Y_b)\geq\tau_{C,U}(b)$ for every $b$ with $\tau_{C,U}(b)>0$. These fiber inequalities are among the trace inequalities, and conversely domination by $v$ implies all trace inequalities on modifications, by the exact formula just proved. Thus the trace competitors for $\Omega$ are exactly the competitors for the ordinary envelope of $v$, and $\mathcal T_\Omega(u;Y)=P_\Omega[v]$.

For $\psi=P_\Omega[v]_{\mathcal I}$, the Lelong numbers along the fibers are those of $v$, since $\psi$ and $v$ have the same multiplier ideals, and the same domination lemma gives $\psi\preceq v$. Hence
\[
 v-\sup_Yv\leq P_\Omega[v]\leq\psi\preceq v.
\]
Since $\psi\leq0$, it is a competitor for the ordinary envelope of $v$, so $\psi\leq P_\Omega[v]$. This proves \eqref{eq:curve-fibration-trace-envelope}.

The potential $v-\sup_Yv$ is a competitor for $P_{\Omega/2}(P_\Omega[v])$. The contact estimate for envelope fixed points, \cref{prop:fixed-point-finite}, gives finite mass, and \cref{lem:smaller-background-mass-lower} gives strict positivity. The two envelope properties give the model condition \eqref{eq:als-model-condition}, so \cref{thm:als-prescribed-lp} gives the stated equations.

Two Green constructions on $C$ differ by a smooth function, by the Poisson argument in the proof of \cref{thm:curve-trace-realization}, so their pullbacks have the same ordinary type. Every background in the statement gives an envelope ordinarily equivalent to this pullback. This proves the asserted independence, at the level of ordinary types rather than of normalized potentials.
\end{proof}

The same construction respects sums and multiplication by positive real numbers at the level of ordinary types. Indeed, the trace measures on the curve add by \cref{thm:trace-data-additivity}, their Green representatives add for a fixed kernel, and pullback preserves sums. Each envelope in a sufficiently large background represents the corresponding ordinary type.
